\hnheader[Zu einfach $\Rightarrow$ Bias (Underfitting). Zu flexibel $\Rightarrow$ Varianz (Overfitting).][Test-MSE $=\sigma^2+\mathrm{Bias}^2+\mathrm{Varianz}$]{Bias \&}{Varianz}{Generalisierung, Underfitting, Overfitting \& Double Descent}
\hnsrc{main\_notes.pdf (Generalization, Kap.\ 8), notes/cs229-notes4.pdf, notes/error-analysis.pdf}

\begin{hngrid}
%
\begin{hnp}[raster multicolumn=2,acc=hnorange]{1}{Intuition}
\textbf{Bias}: Fehler, der bleibt, selbst mit \emph{unendlich} vielen Daten -- die Hypothesenklasse ist zu schwach (z.\,B.\ Gerade statt Parabel).\\
\textbf{Varianz}: wie stark $h_S$ zwischen verschiedenen Trainingssets \emph{gleicher Größe} schwankt -- das Modell lernt Rauschen mit.
\end{hnp}
%
\begin{hnp}[raster multicolumn=2,acc=hnpurple]{2}{Zerlegung (Regression)}
\[ \begin{aligned}\E\bigl[(y-h_S(x))^2\bigr]&=\underbrace{\sigma^2}_{\text{Rauschen}}+\underbrace{(h^*-\bar h)^2}_{\mathrm{Bias}^2}\\&\quad+\underbrace{\E[(\bar h-h_S)^2]}_{\text{Varianz}}\end{aligned} \]
\hnleg{$y$ Beobachtung, $h_S$ auf Trainingsmenge $S$ gelerntes Modell, $h^*$ wahre Funktion, $\sigma^2$ Rauschvarianz, $\E_S$ Erwartung über Trainingsmengen.}
$\bar h=\E_S[h_S]$ ist das ``mittlere Modell'' (Zentrum der Punktwolke rechts).
\end{hnp}
%
\begin{hnp}[raster multicolumn=2,acc=hnteal]{3}{Parameterraum: Zielscheiben}
\centering
\begin{tikzpicture}[scale=.72]
  \foreach \ox/\oy/\sx/\sy/\sp in {0/0/0/0/.15, 3/0/0/0/.7, 0/-3/.75/.55/.15, 3/-3/.75/.55/.7}{
    \begin{scope}[shift={(\ox,\oy)}]
      \draw[gray!50] (-1.2,-1.2) rectangle (1.2,1.2);
      \foreach \a/\b in {.3/.4,-.5/.2,.4/-.6,-.3/-.5,.7/.1,-.2/.7}{
        \fill[red!70!black] ({\sx+\sp*\a},{\sy+\sp*\b}) circle(.09);}
      \fill[hnblue!70!black] (0,0) circle(.1);
    \end{scope}}
  \node[font=\tiny] at (0,1.55) {Bias$\downarrow$ Var$\downarrow$};
  \node[font=\tiny] at (3,1.55) {Bias$\downarrow$ Var$\uparrow$};
  \node[font=\tiny] at (0,-1.45) {Bias$\uparrow$ Var$\downarrow$};
  \node[font=\tiny] at (3,-1.45) {Bias$\uparrow$ Var$\uparrow$};
\end{tikzpicture}\\[-1pt]
{\tiny blau: $\theta^*$, rot: $\hat\theta_S$ verschiedener Datensätze}
\end{hnp}
%
\begin{hnp}[raster multicolumn=3,acc=hnnavy]{4}{Modellkomplexität}
\centering
\begin{tikzpicture}[scale=1.3,>=Stealth]
  \draw[->,hnnavy] (0,0)--(5.2,0) node[right,font=\tiny]{Komplexität};
  \draw[->,hnnavy] (0,0)--(0,2.6) node[above,font=\tiny]{Fehler};
  \draw[hnrose,thick,domain=.2:5,samples=40] plot(\x,{2.2*exp(-.9*\x)+.05});
  \draw[hnpurple,thick,domain=.2:5,samples=40] plot(\x,{.05+.02*exp(.95*\x)});
  \draw[hnorange,very thick,domain=.2:5,samples=40] plot(\x,{.3+2.2*exp(-.9*\x)+.02*exp(.95*\x)});
  \draw[dashed,gray] (2.05,0)--(2.05,2.4);
  \node[font=\tiny,hnrose] at (.9,1.5) {Bias$^2$};
  \node[font=\tiny,hnpurple] at (4.4,.65) {Varianz};
  \node[font=\tiny,hnorange] at (4.2,2.2) {Test-Fehler};
  \node[font=\tiny] at (2.05,-.25) {Optimum};
\end{tikzpicture}
\end{hnp}
%
\begin{hnp}[raster multicolumn=3,acc=hngreen]{5}{Double Descent}
\centering
\begin{tikzpicture}[scale=1.3,>=Stealth]
  \draw[->,hnnavy] (0,0)--(5.2,0) node[right,font=\tiny]{\# Parameter $d$};
  \draw[->,hnnavy] (0,0)--(0,2.6) node[above,font=\tiny]{Test-Fehler};
  \draw[hnorange,very thick,domain=.1:2.3,samples=30] plot(\x,{1.0+.4*exp(-2*\x)+.25*\x+min(1.3,(\x>1.9)*6*(\x-1.9))});
  \draw[hnorange,very thick,domain=2.4:5,samples=30] plot(\x,{2.2*exp(-.9*(\x-2.4))+.55});
  \draw[dashed,gray] (2.4,0)--(2.4,2.4);
  \node[font=\tiny] at (2.4,-.25) {$n\approx d$};
  \node[font=\tiny,hnnavy] at (1.0,2.0) {klassisch};
  \node[font=\tiny,hnnavy] at (4.1,2.0) {überparametrisiert};
\end{tikzpicture}\\
Im überparametrisierten Regime wählt Gradient Descent implizit die \emph{Minimum-Norm-Lösung} -- oft gut.
\end{hnp}
%
\begin{hnex}[raster multicolumn=3]{Beispiel: Polynomgrad wählen (illustrative Zahlen)}
Daten aus einer Parabel + Rauschen.\\
\renewcommand{\arraystretch}{1.2}
\begin{tabular}{@{}c|ccc@{}}
Grad&1&2&5\\\hline
Train-MSE&0.30&0.05&0.00\\
Test-MSE&0.32&0.06&0.45\\
\end{tabular}\\
Grad 1: beide hoch $\Rightarrow$ \hlp{Bias}. Grad 5: Train $\approx0$, Test hoch $\Rightarrow$ \hlp{Varianz}. Grad 2: bester Kompromiss.
\end{hnex}
%
\begin{hnp}[raster multicolumn=3,acc=hnrose]{6}{Diagnose-Regeln}
\begin{itemize}
\item \textbf{Hoher Bias}: Trainingsfehler schon hoch, kleine Lücke zum Testfehler; mehr Daten helfen \emph{nicht}.
\item \textbf{Hohe Varianz}: große Lücke Train/Test; Testfehler sinkt noch mit mehr Daten.
\item Irreduzibles Rauschen $\sigma^2$ kann man nie eliminieren.
\end{itemize}
\end{hnp}
%
\begin{hnp}[raster multicolumn=2,acc=hnred]{7}{Gegen Bias}
\begin{itemize}
\item komplexeres Modell / mehr Features
\item weniger Regularisierung
\item länger trainieren
\end{itemize}
\end{hnp}
%
\begin{hnp}[raster multicolumn=2,acc=hngreen]{8}{Gegen Varianz}
\begin{itemize}
\item mehr Trainingsdaten
\item Regularisierung, weniger Features
\item Bagging / Ensembles
\end{itemize}
\end{hnp}
%
\begin{hnp}[raster multicolumn=2,acc=hnorange]{9}{Key Takeaways}
\begin{hnchk}
\item $\sigma^2+\mathrm{Bias}^2+\mathrm{Var}$
\item Bias: Klasse zu schwach
\item Varianz: Datensatz-Empfindlichkeit
\item Double Descent bei $n\approx d$
\end{hnchk}
\end{hnp}
%
\begin{hnp}[raster multicolumn=6,acc=hnpurple,colback=hnlilac!30]{?}{Selbsttest: kannst du das erklären?}
\hnquiz{Zerlegung des Test-MSE?}{$\sigma^2$ (Rauschen) $+\ \mathrm{Bias}^2+$ Varianz.}{Woran erkennt man hohen Bias?}{Trainingsfehler schon hoch, kleine Lücke zum Testfehler.}{Was ist Double Descent?}{Testfehler steigt bei $n\approx d$ und fällt im überparametrisierten Regime erneut.}
\end{hnp}
%
\begin{hnbanner}[raster multicolumn=6]
„Ein Modell muss so komplex sein wie die Daten -- und keinen Deut komplexer.“
\end{hnbanner}
\end{hngrid}
